\section{Hierarchical Gaussian calculus}
\label{sec:foundations}

The size comparison uses two kinds of calculation: expectations under
hierarchical probability laws, and signed sums of such expectations produced
by covariance differentiation. We define these separately. The resulting
finite-order cavity and size expansions express pressure changes in terms
of overlap errors relative to a one-site scalar model.

\subsection{Normalization and positive tree laws}

We refer to \cite{PE} as PE; references explicitly marked PE concern that
paper. Fix \(T=\beta^2>1\). As in \eqref{eq:completed-free-energy}, we
add an independent spin-constant Gaussian of variance \(T/2\) and use the
uniform probability prior on spins. The resulting matrix covariance is
\(mTQ^2/2\), where \(Q\) is the overlap defined below. The added Gaussian
has bounded moments, and replacing the counting prior by the probability
prior shifts only the mean.

We work in dyads \(N\le m\lesssim 2N\) (with small enlargements), \(m\) the integer number of sites, and put \(x=N^{-1/6}\). All exponents below, including any order of a Taylor expansion, are fixed as \(N\to\infty\); constants can depend on these choices and on \(T\). We use ``with subpower loss'' or \(\lesssim_\circ\) to allow \(C_\epsilon x^{-\epsilon}\) for \emph{every} sufficiently small fixed \(\epsilon>0\). ``Power small'' means \(O(x^c)\) for some fixed \(c>0\); for assertions in every fixed \(L^p\) a power saving not indexed by \(p\) means the exponent is independent of the fixed finite moment order (constant and threshold need not be). Power slack choices in estimates with an indicated hierarchy of exponents are made before using subpower losses.

Estimates in this dyadic notation apply with \(N/2\le m\le3N\); changing comparison scales by bounded factors from this convention is immaterial. Several pieces have local variables (including error majorants and truncation orders) redefined in each calculation. In particular the path martingale \(X_v\) below is unrelated to the notation for the standardized free energy, and in Section~\ref{sec:budgets} the letter \(X\) without a path subscript is also used for a scale ratio.

Fix a terminal mass \(a\in[1,4]\). The coordinate \(v\in(0,a)\)
is the replica mass, and the deterministic nonnegative nondecreasing
functions \(b,h\) are the matrix and spin-field clocks. Jumps are allowed.
Their cumulative Gaussian covariance at mass \(v\) is
\[
 m(b(v)Q^2/2+h(v)Q),\qquad
 Q=m^{-1}\sum_{i\le m}\sigma_i\sigma'_i.
\]
For step clocks, construct independent Gaussian increments with these
covariance increases, including any initial increment at mass zero.
If \(H\) is their sum, the terminal continuation is
\(a^{-1}\log\int e^{aH}\), with the integral taken against the uniform
probability prior. Compute earlier continuations backwards: for an
increment at mass \(v>0\), set
\[
 F_{\rm before}=v^{-1}\log\mathbf E\exp(vF_{\rm after}),
\]
where this expectation averages only that increment. At mass zero use
ordinary expectation. Divide the resulting deterministic value by \(m\)
and subtract \(a(b(a-)/4+h(a-)/2)\); this defines the self-corrected
pressure \(f_m(b,h)\). Its interval length \(a\) will be implicit.

The same continuations define a forward probability law. Reveal each
Gaussian increment with its original Gaussian density multiplied by
\(\exp(v(F_{\rm after}-F_{\rm before}))\), and then draw the terminal
spin from the Gibbs law of exponent \(a\). The backward recursion makes
each transition integrate to one. On a fixed finite rooted tree, branches
share the increments before their fork and use conditionally independent
forward transitions afterwards. We call this the \emph{positive tree law}
and write \(\nu[P]\) for its expectation of \(P\), including both fields
and terminal spins. The tree topology and clocks are fixed in this
expectation. Bounded monotone clocks are obtained by the fixed-site step
limits from PE stated below.

A specified observation may lie on either side of a jump or at an
intermediate clock value. At a simultaneous jump we use the joint filling
\[
 (b_t,h_t)=(b_-,h_-)+t(\Delta b,\Delta h),\qquad 0\le t\le1,
\]
with mass held fixed, and preserve this filling in step approximations.
Endpoint tree laws and mass-integrated identities do not depend on the
order of revelation within a same-mass jump. The separate matrix and
field contributions to quadratic variation, and observations at
intermediate cuts, do use the prescribed filling.

For a single path let \(y=\sigma/\sqrt m\),
\[
\begin{gathered}
 X_v=\mathbf E_v y,\qquad C_v=\mathbf E_v(y-X_v)(y-X_v)^{\mathsf T},\\
 S_v=\mathbf E|X_v|^2,\qquad
 B_v=\mathbf E\|\mathbf E_v yy^{\mathsf T}\|_{\rm HS}^2,\qquad D_v^2=B_v-S_v^2.
\end{gathered}
\]
Here \(\mathbf E_v\) is conditional expectation given the fields revealed
through the specified cut, and \(\mathbf E\) averages the full single-path
positive law. Thus \(X_v,C_v\) are functions of the random prefix, whereas
\(S_v,B_v,D_v\) are deterministic statistics of that law. Write \(H_v\)
for the Hessian of the continuation with respect to a source coupled to
\(y\) in the downstream Hamiltonian, with the revealed prefix held fixed. Put
\(L_v=y\cdot X_v-|X_v|^2/2\) and
\(W_v=(y-X_v)^{\mathsf T}H_v(y-X_v)\). These last two observables also
retain the terminal spin.

\subsection{Pinned moments and projection}

The next estimates distinguish a variation of the deterministic clocks
from movement along one fixed revealed path. In the pressure formula,
\(db(v),dh(v)\) are clock variations in an external parameter. In
the susceptibility inequality, \(dS_v\) and \(dh\) are instead
Stieltjes increments along the fixed path, including its prescribed
filling at a jump. The source Hessian \(H_v\) always keeps the
revealed prefix fixed.

\begin{lemma}[Pressure, source, and single-path estimates]
\label{found:estimates}
For the positive laws just defined, the following identities and
inequalities hold at every fixed size. The moment bounds in
\eqref{eq:A1}--\eqref{eq:A2} hold for every fixed real \(2\le p<\infty\),
with deterministic \(k\in L^2(0,a)\) and zero deterministic external
field.
\[
\begin{split}
 df_m&=-\tfrac14\int B_v\,db(v)\,dv-\tfrac12\int S_v\,dh(v)\,dv,\\
 vC_v\preceq H_v\preceq a C_v,\qquad
 d S_v&\ge m\,\mathbf E\operatorname{tr}H_v^2\,dh,\qquad
 \mathbf EW_v=\mathbf E\operatorname{tr}C_v H_v,\\
 \left\|\int k_v(L_v-\mathbf EL_v)dv\right\|_p^2&\le c_p\left\|\int k_v^2 W_v\,dv\right\|_{p/2}.
\end{split}                                                    \tag{A1}\label{eq:A1}
\]
The last line is a moment estimate under the full single-path law, with
deterministic weights \(k_v\) and zero deterministic external field.
More generally, for
polynomials \(P_v\) of fixed degree with bounded deterministic coefficients,
\[
 \left\|\int k_v(P_v(L_v)-\mathbf E P_v(L_v))dv\right\|_p^2
 \le c_p\left\|\int k_v^2 P'_v(L_v)^2 W_v\,dv\right\|_{p/2}.
                                                               \tag{A2}\label{eq:A2}
\]

For the projection bound only, assume that the total clocks are at most
a fixed power of \(\log N\), and that the positive mass and field
length are bounded below by inverse powers of \(N\). If a branch
has traversed a pure or mixed clock interval with field length
\(\Delta h>0\) after separating from the data determining an anchor
vector \(V\), and its mass is at least \(v>0\) throughout that
interval, then for every fixed even \(p\ge2\),
\[
 \|V\cdot(y-X_{\rm end})\|_p
 \lesssim_p \frac{\log^2 N}{v\sqrt{m\Delta h}}\,\|V\|_p                       \tag{A3}\label{eq:A3}
\]

where the vector norm inside \(L^p\) is Euclidean. The branch must
retain its normalized future kernel conditional on the prefix and outside
data; in particular, the bound applies to conditionally independent
resamplings.
\end{lemma}

\begin{proof}
Equations \eqref{eq:A1} and \eqref{eq:A3} are the pressure, source,
pinned-moment, and projection estimates of \cite[Lemmas~2.1--2.3
and~3.5, equations~(1) and~(13)]{PE}. To match the terminal-mass-one
normalization there, put
\[
 u=v/a,\qquad \widetilde b(u)=a^2b(au),\qquad
 \widetilde h(u)=a^2h(au),\qquad \widetilde F=aF.
\]
The conditional spin means and covariances are unchanged, while
\(H_{au}=a\widetilde H_u\). The projection denominator is unchanged
as well, since \((v/a)\sqrt{m a^2\Delta h}=v\sqrt{m\Delta h}\).
Appendix~\ref{sec:companion-interface} records the pressure and moment
conversions. The fixed-size path-limit result is PE, Proposition~2.14.

For \eqref{eq:A2}, we include the polynomial argument of PE,
Appendix, paragraph ``Source and pressure calculus; the scope of pinned
moments.'' Pin the terminal spin and write the Gaussian increments in
terms of standard Gaussian drivers, as in the proof of PE, Lemma~2.3. Let \(\mathsf A_v\) be the Hessian of the continuation
in these drivers and \(\mathsf B_v\) its mixed Hessian with the
normalized spin source. After the transition densities telescope, the
negative logarithm of the pinned driver density is the Gaussian
quadratic potential minus \(aH_\sigma\), which is linear, plus
\(\int_0^a F_v\,dv\), where \(F_v\) is the continuation at the cut
\(v\). Consequently
\[
 \begin{pmatrix}\mathsf A_v&\mathsf B_v\\
 \mathsf B_v^{\mathsf T}&H_v\end{pmatrix}\succeq0,
 \qquad
 \mathsf Q=I+\int_0^a\mathsf A_v\,dv,
\]
where \(\mathsf Q\) is the Hessian of the negative logarithm of the
pinned driver density. The driver gradient of \(L_v\) is
\(g_v=\mathsf B_v(y-X_v)\). Thus, for every driver direction \(z\),
\( |z\cdot g_v|^2\le(z^{\mathsf T}\mathsf A_vz)W_v\).
For the polynomial observable set
\[
 G_P=\int_0^a k_vP'_v(L_v)g_v\,dv,
 \qquad R_P=\int_0^a k_v^2P'_v(L_v)^2W_v\,dv.
\]
Cauchy--Schwarz in \(v\) gives
\( |z\cdot G_P|^2\le(z^{\mathsf T}\mathsf Qz)R_P\), hence
\(G_P^{\mathsf T}\mathsf Q^{-1}G_P\le R_P\).
The inverse-Hessian Brascamp--Lieb inequality~\cite[Theorem~4.1]{BL1976}
therefore gives \eqref{eq:A2} for \(p=2\). Applying the same variance
inequality to the \(p/2\)-th absolute power of the centered integral,
as in PE, Lemma~2.3, gives the assertion for every fixed \(p>2\).

The zero-field spin gauge transformations act transitively on terminal
spins and preserve both \(L_v\) and \(W_v\). Because the coefficients of \(P_v\) and the weights \(k_v\)
are deterministic, they also preserve the polynomial integral and
\(R_P\). Their conditional laws are consequently independent of the
pinned spin, including the mean used for centering. Averaging over that
spin yields the full single-path estimate. Bounded-weight approximation
and then the fixed-\(m\) path limit give the stated range of clocks
and weights.

Clocks are bounded for each \(N\), possibly by a power of \(\log N\); \eqref{eq:A1} is clock independent. For \eqref{eq:A3}, we apply PE, equation~(13), only in this bounded or polylogarithmic clock regime; together with the polynomial mass and field-length bounds, this gives the polynomial coefficient-ratio bounds needed in its moment iteration, as detailed in Appendix~\ref{sec:companion-interface}. Our estimates can always be passed from affine finite-step interpolations, or from inserted fixed cuts, by the fixed-\(m\) limit assertion, keeping chosen revelation sides and the joint filling at observed intermediate jump cuts, and using the corresponding integrated inequalities. In particular no discretization \emph{rate} in \(m\) is used.

The pinned inequalities in this lemma concern the full single-path law. When a tree contains additional branches, we first marginalize the branches not used by the test. Random factors retained on other branches are subsequently handled by unconditional H\"older inequalities; they are not made into random weights in \eqref{eq:A2}.
\end{proof}

\subsection{Signed allocation of replica labels}

Covariance derivatives introduce additional replica labels. For a simple
example, start with one observed leaf above a cut \(v\), and divide its
edge at a level \(u\in(v,a)\). An extra label may coincide with that
leaf, with coefficient \(a\), or split from it along either edge
interval, with coefficient minus the interval length. On a constant
summand the total is \(a-(a-u)-(u-v)=v\). In a larger tree, attachment
to a fork also contributes an atom. We sum all these placements by the
signed rule of PE, Section~2,
Proposition~2.5, equation~(2), and Lemmas~2.6, 2.8, and~2.9. A new
label coinciding with an observed leaf has coefficient \(a\); a new
split on an occupied edge has coefficient \(-dv\); attachment at an
existing mass-\(v\) fork with \(d\) children has coefficient
\(-(d-1)v\). These coefficients multiply expectations under the
positive laws just defined. They are not probabilities.

The row-sum rule sums one free label whose summand is constant over all
its attachments in a block above a cut \(v\). The result is that constant
times \(v\), as in the example. A row whose summand depends
on its attachment cannot be collapsed this way. Empty levels can be
inserted continuously without changing the calculation. Different
infinitesimal fork variables specify different tree vertices even when
the covariance clocks have the same value at those vertices.

A sum \emph{per root mass}, denoted by \((\cdot)_\#\), begins without
specified old labels. Its first label is an ordinary observation with
weight 1, and all subsequent labels use the signed rule. This is the
limit of normalized summation as the incoming mass tends to zero. Sums
without \(\#\) extend whatever old labels are already present.
Products of symbolic slot polynomials concatenate all their labels;
a coefficient involving only some labels is evaluated on the restriction
of the final topology to those labels. The labeled product rules in PE,
Proposition~2.5 make this convention associative, commutative, and
independent of allocation order. In a per-root product we concatenate first and apply the root
normalization only once.

Constant diagonal covariance changes contribute only to the self
correction. Thus pressure insertions after that correction use ordered
off-diagonal pairs and a factor \(1/2\) per insertion. The first
insertion gives \eqref{eq:A1}; for subsequent insertions the constant
diagonal terms cancel by the row-sum rule in PE, equation~(2). The same
rule applies to any affine positive semidefinite interpolation between
finite feature models with constant diagonal and fixed \(a\), after
subtracting \(a/2\) times the terminal diagonal. For total pressures,
the inserted covariance change is the total covariance change, rather
than its per-site value. Deterministic additive terms common to the two
ensembles cancel when we take their pressure difference.

Enlarging a positive-law tree preserves its original marginals. For any
fixed number of inserted labels, the total variation of the signed
allocation coefficients is uniformly bounded. Taking absolute values
in placement order leaves finitely many integrals over new fork masses,
existing-fork atoms with bounded multiplicities times their masses,
and terminal coincidences. We may therefore reorder the exact signed
sum with its topology restrictions first, and then estimate with this
absolute allocation measure. This order of operations is essential
when a later argument uses cancellation.

\subsection{Scalar references and Laplace tests}

The signed calculus will compare the matrix model with independent
one-site scalar systems. For the reference model take one spin and only a deterministic spin clock
\(K(v)\). Its continuation can be parametrized by cumulative variance
\(t\) instead of mass. The mass already reached by time \(t\) is
\(\alpha(t)=|\{v\in(0,a):K(v)\le t\}|\). Write \(V(t,z)\) for
the scalar continuation at field \(z\), and \(Z_t\) for its field under
the forward positive law. They satisfy
\[
 V_t=-\tfrac12(V_{zz}+\alpha V_z^2),\qquad dZ_t=\alpha V_z(t,Z_t)dt+dB_t ,
 \qquad V_{\rm term}=a^{-1}\log\cosh(a z).
\]
Set \(\Gamma(t)=\mathbf E V_z(t,Z_t)^2\), with zero starting field
unless otherwise stated. A terminal mass-\(a\) extension in scalar time
changes earlier continuation values only by a spatial constant. In
particular it leaves their derivatives in \(z\) unchanged. On a fixed
tree let \(\tau_i\) be the scalar spin at leaf \(i\), and let
\(v_{ij}\) be the mass where distinct leaves \(i,j\) separate.
Their positive-law moment is
\(\nu_K[\tau_i\tau_j]=\Gamma(K(v_{ij}))\). Common covariance insertions and source derivatives of these moments again use the tree calculus. In particular any fixed spatial derivative of positive order of \(V\) is bounded uniformly (even with long but finite clocks). Below put
\[
 q(v)=\Gamma(K(v)),\qquad h_*(v)=K(v)-b(v)q(v),\qquad d(v)=h(v)-h_*(v).
\]

The physical paths used below have \(b(v)=b(0+)\) and \(h(v)=0\)
for \(v<x\). For \(0<s<x\), move the initial matrix revelation
from mass zero to mass \(s\). This gives the tilted pressure
\(f_m^s=f_m(b\,{\bf1}_{v\ge s},h)\); the untilted pressure is
\(f_m^0=f_m(b,h)\). Put \(\psi_m(s)=f_m^s-f_m^0\). The factor
\({\bf1}_{v_{ij}\ge s}\) records whether two tilted branches share
the matrix revelation. Accordingly define \(Q^*_{ij}=Q_{ij}\) in
the untilted ensemble and
\(Q^*_{ij}={\bf1}_{v_{ij}\ge s}Q_{ij}\) in the tilted ensemble.
The mass-\(s\) power mean then gives
\[
 {\cal K}_m(s)=s\,m\psi_m(s)=\log\mathbf E\exp\{s(G_m-\mathbf E G_m)\},
                                                               \tag{A4}\label{eq:A4}
\]
where \(G_m\) is the random total continuation just after the initial matrix revelation in the untilted path. Each connected subtree with internal splits at \(v\ge s\) in a tilted positive law has its untilted marginal reweighted by \(\exp(sG_m)/\mathbf E\exp(sG_m)\). Distinct such components have independent roots. Products using active branches from more than one component can therefore be bounded by their active-subtree marginal norms without interpreting the forest as sharing an untilted root. Set
\[
 \Delta=Q^*-q,\qquad J=h+b Q^*-K=b\Delta+d
\]
on off-diagonal pairs. Scalar paths used as references are \emph{unshifted} in both ensembles. Tests involving matrix-model overlaps only through \(Q^*\) are invariant under simultaneous sign flip of all spins on a component (with its high ordinary fields negated); consequently at every pair \(ij\), including below \(s\) in the tilted case, \(\nu[P Q_{ij}]=\nu[P Q^*_{ij}]\) for these spin tests.

\subsection{Finite row and size expansions}

We now encode fixed finite Taylor expansions in the row mismatch \(J\).
A \emph{jet} means that the series is retained only through a specified
finite degree; no convergence of an infinite series is being asserted.
In all inserted pairs, the two ends must be different leaves, although
labels from different pairs may coincide. Define the symbolic series
\[
 \Psi_K(J)=\left\langle\exp\{\tfrac12\sum J_{cd}\tau_c\tau_d\}\right\rangle_K,
 \qquad N_{ij,K}(J)=\left\langle\tau_i\tau_j
                \exp\{\tfrac12\sum J_{cd}\tau_c\tau_d\}\right\rangle_K .         \tag{A5}\label{eq:A5}
\]
To read \eqref{eq:A5}, first expand the exponential to the required
finite degree and allocate its pair labels. For each final topology,
the brackets are the positive scalar moment of the labels in that
coefficient. Multiply coefficients on their respective restricted
topologies before carrying out the signed allocation sum. Thus
\(\log\Psi\) has the ordinary joint scalar cumulants at each fixed
topology as its coefficients. Products, differentiation, and inversion
of the series with constant term 1 are all finite-degree polynomial
operations under this convention.

For example, in the product of two
scalar coefficients \(\nu_K[\tau_i\tau_j]\nu_K[\tau_c\tau_d]\), allocate
all four labels on one final topology, then evaluate the first moment
on its \(ij\) restriction and the second on its \(cd\) restriction.
Their product is not the four-spin moment, even when some labels
coincide. The clock derivative and the coefficient linear in the number
of added sites will be expressed through the same pressure symbol,
namely
\[
 {\cal I}=\log\Psi_K(J)-\tfrac14\sum b_{ij}(Q^*_{ij})^2 .
                                                               \tag{A6}\label{eq:A6}
\]
Independent additive deterministic per-root terms in this symbol will always cancel.

The following Taylor rules require a monotone reference \(K\) and
bounded \(h-K\) and \(b\). A remainder with \(k\) insertions is a
finite sum over boundedly many allocated labels, possibly integrated
over interpolation parameters. Its integrand is a bounded factor times
\(k\) specified overlap insertions on off-diagonal pairs. We estimate
such terms using the absolute allocation measure and a positive
probability law.

The intermediate law may still contain one interpolated scalar row.
For a nonnegative spin test of the other rows, PE, Lemma~2.9 compares
this law to the actual law up to a constant: its Gr\"onwall bound uses
the bounded total covariance change of that row. Deleted sites are
omitted from an overlap without changing its denominator \(m\).
Restoring those sites changes a bounded Lipschitz polynomial of fixed
degree by \(O(m^{-1})\). Any deterministic topology coefficient or
weight multiplying the test also multiplies this error bound. Throughout
this interpolation the spin test is fixed, so it contributes no
additional derivative.

\begin{lemma}[Finite cavity quotient]
\label{found:cavity-quotient}
For overlap-polynomial \(P\) invariant as after \eqref{eq:A4}, on any fixed finite-label tree including \(ij\), one has through order \(k-1\)
\[
 \nu[P Q^*_{ij}]=\nu[P(N_{ij,K}/\Psi_K)_{<k}]
          +{\rm Rem}_k+O(m^{-1}),                                      \tag{A7}\label{eq:A7}
\]
where \({\rm Rem}_k\) has at least \(k\) insertions \(J\), times \(P\).
\end{lemma}

\begin{proof}
By exchangeability, the observed overlap may be represented using one
chosen site. Delete that site's contribution from the other tested
overlaps, retaining the denominator \(m\), and interpolate its Gaussian
kernels to those of an independent one-site scalar model with clock
\(K\). All other covariance kernels remain fixed. Both endpoints are
positive semidefinite, as is their affine interpolation. In the expansion
back from the independent endpoint, the covariance insertion is
\(J_{cd}\tau_c\tau_d\), where \(J\) uses the remaining-row overlap
and the mask of the actual matrix clock. The omitted terms are
\(O(m^{-1})\) and constant diagonals. The row covariance change is
bounded, so Taylor's integral remainder has the form described above.

At the independent endpoint, averaging the removed spin gives
\(\Psi_K\) when it is unobserved, and \(N_{ij,K}\) when its product
on leaves \(i,j\) is observed. Expanding the former against the
truncated quotient \(N_{ij,K}/\Psi_K\) cancels its factor \(\Psi_K\)
to the retained order. Allocation order symmetry identifies the terms
on each final topology and gives \eqref{eq:A7}. Products here multiply
separately averaged scalar coefficients; they do not couple their
scalar spins. In particular the linear cavity kernel on a pair \(cd\) is
\[
 {\bf M}_{ij,cd}=\tfrac12
 (\nu_K[\tau_i\tau_j\tau_c\tau_d]-\nu_K[\tau_i\tau_j]\nu_K[\tau_c\tau_d]).
                                                                    \tag{A8}\label{eq:A8}
\]
This proof also gives the unestimated remainder representation using one interacting scalar row (the removed site) along the affine interpolation, if symmetries must be used before taking absolute values.

\end{proof}

\begin{lemma}[Change of scalar reference]
\label{found:reference-change}
Apply \eqref{eq:A7} at a monotone reference \(K'\), then Taylor-expand
its deterministic scalar coefficients from \(K\) to \(K'\). At any
consistent total truncation, the result is the quotient with reference
\(K\) and combined argument
\(J=(h+bQ^*-K')+(K'-K)\).
\end{lemma}

\begin{proof}
PE covariance differentiation, equation~(2), represents a derivative
of a scalar coefficient by a new pair insertion
\((K'-K)_{cd}\tau_c\tau_d/2\). Allocation order symmetry permits us
to combine these with the insertions already present in the quotient.
The apparent extra factors \(\Psi_K(K'-K)\) have no effect on an
observed topology whose spin test is independent of the new labels:
after allocating these labels last, every positive-degree term is a
derivative of the unit scalar moment and sums to zero. For a standalone
per-root sum the same terms can remain, but they are deterministic.
Taylor's integral remainder has bounded scalar coefficients and powers
of \(K'-K\); it can also retain ordinary polynomial factors of the
original argument \(h+bQ^*-K'\).
\end{proof}

\begin{lemma}[Differentiated row identity]
\label{found:row-derivative}
Let a dot vary \(b,h,K,q\) at fixed interval length \(a\), and let
\({\cal D}\) differentiate \eqref{eq:A6} including its scalar coefficients,
with the formal rule \({\cal D}\Delta=-e\Delta\) for any fixed \(e\).
Thus \({\cal D}Q^*=\dot q-e\Delta\); this rule acts on the symbolic
polynomial and does not differentiate a random spin sample. Up to
deterministic terms common to the two ensembles and Taylor remainders,
\[
 \dot f_m^{\,s}\ \equiv\ (\nu[{\cal D}{\cal I}])_\#.                   \tag{A9}\label{eq:A9}
\]
The same formula applies with superscript \(s=0\). The error and its differentiated weights are as described in the proof below.
\end{lemma}

\begin{proof}
Differentiating a scalar coefficient inserts
\(\dot K_{ij}\tau_i\tau_j/2\). This cancels the \(-\dot K\)
part of \({\cal D}J\); terms with no previously attached labels
are deterministic. Thus the logarithm contributes
\[
 \frac12\sum\frac{N_{ij,K}}{\Psi_K}
       (\dot h+\dot b Q^*+b{\cal D}Q^*)_{ij}.
\]
Apply \eqref{eq:A7} to replace the quotient by \(Q^*\), with its
stated error. The term \(bQ^*{\cal D}Q^*/2\) then cancels the
corresponding derivative of the quadratic term in \eqref{eq:A6}.
In particular the choice of \(e\) disappears. The terms left are
\(\frac12\sum\dot h Q^*+\frac14\sum\dot b(Q^*)^2\), precisely
the off-diagonal pressure insertions in \eqref{eq:A1}.

For the error, all tests have bounded degree and at most one
differentiated weight. If \(\log\Psi\) is truncated after degree
\(k\), the residuals in the difference of the two pressure derivatives
have at least \(k\) insertion factors. They also have one weight at
a split, possibly on an additional pair, bounded by a constant times
\[
 1+|\dot h|+|\dot b|+|\dot K|+|\dot q|;
\]
the 1 is omitted when \(e=0\). The \(O(m^{-1})\) errors carry the
same weights. This follows from the product rule, since differentiating
one insertion can remove that factor, and from \eqref{eq:A7}. If a
second reference \(K'\) is used, expand the quotient in this gradient
first, then recombine by Lemma~\ref{found:reference-change}. This does
not require differentiating \(K'\).

There is one truncation convention to check. Differentiating an
insertion lowers its insertion degree by one, whereas differentiating
a scalar coefficient does not. Through degree \(k\) in the old
\(J\) factors, the coefficient part of \({\cal D}\Psi\) therefore
inserts \(\frac12\sum(N_{ij,K})_{\le k}\dot K_{ij}\), formally
including its degree-zero term. In a product with old labels independent
of the new placements, that term sums to zero when the new labels are
allocated last: it differentiates a unit scalar observation. Standalone
per root, its discrepancy is common and deterministic. Thus coefficient
differentiation of the logarithm uses the reciprocal jet through degree
\(k\), while differentiation of an insertion uses the quotient through
degree \(k-1\). All these cancellations take place in the unrestricted
signed sums, before imposing split-scale restrictions for estimates.

\end{proof}

\begin{lemma}[Finite increments in the number of sites]
\label{found:size-increment}
At fixed clocks and \(s\), assume that \(bq^2\) is nondecreasing. For every fixed positive integer \(j\), the change from \(m\psi_m\) to \((m+j)\psi_{m+j}\), expanded through total order \(k\), is a polynomial in \(j\) of degree at most \(k\) with zero constant term. Its linear coefficient is the difference of the per-root expectations of \eqref{eq:A6}, with \(\log\Psi\) truncated after degree \(k\). The remainder has the total-pressure bounds given below.
\end{lemma}

\begin{proof}
For size changes at fixed clocks and \(s\), compare \((m+j)\psi_{m+j}\) to \(m\psi_m\), with any fixed positive integer \(j\). Interpolate the enlarged kernel from that of \(m\) old spins plus \(j\) independent scalar rows at \(K\), to its actual value plus independent, spin-constant cumulative covariance \(j b q^2/2\) (assume \(b q^2\) nondecreasing). This addition gives the same pressure shift for both tests. The off-diagonal insertion before the factor \(1/2\) is
\[
 J\sum_{h'=1}^j\tau^{h'}\tau^{h'\prime}
       -\tfrac j2 b(Q^{*2}-q^2)+O_j(m^{-1}).
\]
Indeed the normalization of the matrix feature changes from \(m\) to \(m+j\) while field terms are additive. This difference is bounded. Taylor through order \(k\) in the total pressure now gives a polynomial in \(j\) of degree \(\le k\), constant term zero; its coefficient linear in \(j\) is the compared per-root expectation of \eqref{eq:A6}, \(\log\Psi\) truncated after degree \(k\). For the linear coefficient, group the labeled insertions by their
assigned scalar row. The scalar rows are independent on each fixed
topology, so their series is \(\Psi_K^j\). If \(C\) denotes the
matrix counterterm \(\frac14\sum b(Q^{*2}-q^2)\), the combined series
is \(\Psi_K^j e^{-jC}\). Its coefficient of \(j\) is
\[
 [j]\big(\Psi_K^j e^{-jC}\big)=\log\Psi_K-C.
\]
The deterministic \(q^2\) term cancels between the two tests, leaving
\eqref{eq:A6}. This identity is taken at the same finite total Taylor
order as the covariance expansion. Remainders after taking centered differences can be bounded by \(O(m^{-1})\) and products of \(k+1\) factors \(|\Delta|+|J|\) at the insertions, under the actual \(m\)-spin laws by the same Gr\"onwall comparison. No \((m+j)\) factor is to be put in front of this remainder.
\end{proof}

These rules are identities and bounds for fixed finite Taylor orders, not infinite replica expansions or analytic continuations in a replica count. Differentiation along absolutely continuous paths in \eqref{eq:A9} uses pointwise (or integrated) chain rules, with deterministic dotted clocks absolutely integrable including their bounded-density fork averages.


The clock derivative and the coefficient linear in the number of added
sites are thus expressed by the same finite pressure symbol
\eqref{eq:A6}. This shared expression permits the size comparison to
combine their signed coefficients before estimating the remainders.
